\section{Bipartite completion of old/new spectral disjointness}
\label{sec:bipartite-spectral}
The remaining bipartite case requires additional mixed entropy phases
and a two-dimensional excess threshold at zero. We prove the exact
companion moments needed for this case before completing the transfer.

\begin{theorem}[All simple regular bases]\label{lift:thm:all-base}
Let $B$ be the adjacency matrix of any fixed connected simple
$d$-regular base, $d\ge3$, and let $A_N$ be its independent uniform
random $N$-lift. With $U_N$ the real fiber-constant space and
$W_N=U_N^\perp$,
\[
 \Prob\{\spec(A_N|_{W_N})\cap\spec(B)\ne\varnothing\}
 \longrightarrow0
\]
through all positive lift sizes. Every old eigenvalue retains exactly
its base multiplicity, including irrational and repeated old modes.
\end{theorem}

The nonbipartite case is Theorem~\ref{lift:thm:old-new}.
For the rest of this section assume the base is bipartite. Regularity
gives equally many base vertices on its two sides, say $m$ each, and
\begin{equation}\label{eq:bip-square-blocks}
 B=\begin{pmatrix}0&T\\T^T&0\end{pmatrix},\qquad
 A_N=\begin{pmatrix}0&X\\X^T&0\end{pmatrix},
 \qquad T\in\Z^{m\times m},\quad X\in\Z^{mN\times mN}.
\end{equation}
These square blocks are essential to the zero-target argument.
Write the two old parts as $\mathcal L,\mathcal R$.
Fix an odd prime $p$ and a monic irreducible $g\in\Fp[t]$ of
degree $r$. For a companion $R_g$, the evaluation isomorphism
\eqref{eq:lift-hom-kernel} gives
\[
 Z_g(M):=|\{Y:MY=YR_g\}|=|\ker g(M)|,
 \qquad h=\dim_{\Fp}\ker g(B).
\]
This is geometric kernel dimension, even when a modular matrix has
Jordan blocks. All limits below fix $B,p,g$ before $N\to\infty$.
We use the phase, relative local limit and weighted sparse estimates
of Section~\ref{sec:frame-moment-proof}, using the positive-weight
phase specialization in Corollary~\ref{ch:cor:positive-phases}. The new compatible indices are proved below.

\begin{theorem}[Bipartite companion moments]\label{bip:thm:moments}
For every positive lift-size sequence,
\begin{equation}\label{eq:bip-moments}
 \E Z_g(A_N)=c_g p^h+o(1),\qquad
 c_g=\begin{cases}
 p+3,&g=t,\\
 2,&g\ne t\text{ and }g\text{ is not even},\\
 2(p^s+1),&g\text{ is even of degree }2s.
 \end{cases}
\end{equation}
Here even means $g(-t)=g(t)$. No condition on $p\mid N$ or
$p\mid d$ is imposed.
\end{theorem}

\subsection{The full-uniform coefficient and its extra forms}
Use row labels $a\in V=\Fp^{1\times r}$, let $q=p^r$, and
temporarily allow any target $R$. The maximal legal-star congruences
are \eqref{eq:lift-congruences}; before matching compatibility their
index is $p^{krq}$, where $k=2m$. The compatible integer lattice
has the free root-margin and edge-interior parametrization used in
the proof of Lemma~\ref{lift:lem:cyclic}. Its reduction modulo $p$
is the whole compatible direction space, with no saturation factor.

A dependency of the congruences has columns $h_i(a)$.
The same rectangular edge transfers force
\[
 h_i(a)=K_i a^T+c_i,\qquad K_j=-K_i^T\quad(i\sim j).
\]
Connectivity now leaves $K_i=M$ on the left and $K_i=-M^T$ on
the right, with no odd-cycle skew constraint. The independent
root-margin tests at $a$ and $-a$ give
\[
 \sum_{j\sim i}c_j=Rc_i,\qquad
 RM\text{ and }RM^T\text{ are skew-symmetric}.
\]
Conversely these conditions annihilate every compatible direction.
The constant columns form the Frobenius-adjoint old intertwiner
space and have dimension $h_B(R)=\dim\{Y:BY=YR\}$. Put
\[
 b(R)=\dim_{\Fp}\{M:RM\text{ and }RM^T
                                      \text{ are skew-symmetric}\}.
\]
The restricted congruence rank is therefore
$\rho=krq-h_B(R)-b(R)$.

This rank gives an exact coefficient. The relative legal/independent
star density ratio on the logarithmic central window is
$p^{krq}(1+o(1))$. The compatible legal sublattice has index $p^\rho$.
Its Gaussian sum has fraction $p^{-\rho}$ of the independent one,
whose exact weighted normalizer is one by
\eqref{eq:lift-two-normalizers}. The mean, covariance and matching
Stirling factors are the same as in the preceding lift calculation;
the positive precision is \eqref{eq:lift-weighted-precision} and
all powers of $N$ cancel. The all-zero legal stars provide a
compatible affine point for every $N$. Translated lattice sums and
the relative local limit estimate therefore retain this fraction
at every residue, and the fixed-neighborhood tails are $o(1)$.
Thus the full-uniform contribution is exactly
\begin{equation}\label{eq:bip-uniform}
 p^{h_B(R)+b(R)}+o(1).
\end{equation}

If $R$ is invertible, set $L=RM$. The two skew conditions are
equivalent to
\begin{equation}\label{eq:bip-form}
 L^T=-L,\qquad R^2L=L(R^T)^2.
\end{equation}
Indeed $M=R^{-1}L$, and skewness of $RM^T=-RL R^{-T}$ is
equivalent, after multiplication by $R$ and $R^T$, to the second
identity. Hence $b(R)$ counts alternating forms for which right
multiplication by $R^2$ is self-adjoint, including degenerate forms.

For an irreducible $g\ne t$, identify this row module with
$E=\Fp[\lambda]$ and put $K=\Fp[\lambda^2]$.
The extension $E/K$ has degree at most two. Its degree is two
exactly when $g$ is even: a nontrivial automorphism sends
$\lambda$ to $-\lambda$, while an even minimal polynomial supplies
that automorphism. The nonzero constant term forces the degree of
an even $g$ to be $2s$, with $|K|=p^s$.
If $g$ is not even, $K=E$, so $R^2$ is cyclic; the cyclic
alternating-form argument forces $b(R)=0$.

In the even case $E$ has dimension two over $K$.
The forms in \eqref{eq:bip-form} are precisely the $K$-balanced
alternating $\Fp$-forms $\beta(a,b)=aLb^T$, where balance means
$\beta(cx,y)=\beta(x,cy)$ for $c\in K$.
They correspond to alternating $K$-bilinear
forms $\Omega$ through $\beta(x,y)=\operatorname{Tr}_{K/\Fp}\Omega(x,y)$.
To justify this correspondence, the trace polynomial
$z+z^p+\cdots+z^{p^{s-1}}$ is nonzero and has degree below $|K|$;
therefore the pairing $\operatorname{Tr}(cz)$ is nondegenerate.
Define $\Omega(x,y)$ uniquely by
\[
 \operatorname{Tr}(c\Omega(x,y))=\beta(cx,y)\quad(c\in K).
\]
Uniqueness and balance give $K$-linearity in each argument, and
skewness plus odd characteristic gives alternation. Conversely the
trace of any alternating $K$-form has these properties. The space
of alternating forms on a two-dimensional $K$-space has dimension
one over $K$, so $b(R)=s$. This proof also covers $p\mid s$ and
never assumes $\operatorname{Tr}(1)\ne0$.
For $g=t$, the scalar target is zero, both skew conditions are
vacuous, and $b(0)=1$.

\subsection{Nonzero targets and every mixed line phase}
The positive-weight bipartite phase classification
gives subspaces $H_L,H_R$ with
$\lambda H_L\subseteq H_R$ and $\lambda H_R\subseteq H_L$.
For nonzero irreducible targets, invertibility implies
\[
 H_R=\lambda H_L,\qquad \lambda^2H_L=H_L.
\]
If $g$ is not even, irreducibility of multiplication by $\lambda^2$
leaves only the point and full-uniform phases. The point family
has $p^h$ members, each contributing $1+o(1)$ by the sparse bound.
Equation~\eqref{eq:bip-uniform} contributes another $p^h+o(1)$.
This proves the non-even case of \eqref{eq:bip-moments}.

If $g$ is even of degree $2s$, the proper nonzero possibilities are
the $K$-lines
\[
 H_L=uK,\qquad H_R=\lambda uK,
 \qquad \#\{H_L\}=p^s+1.
\]
Coordinates $xu$ on the left and $y\lambda u$ on the right make
the two root maps $I$ and $C$, where $C$ is multiplication by
$c=\lambda^2$ on $K$. This $C$ is cyclic over $\Fp$.
Choose any complementary $K$-line $u'K$ and its right partner
$\lambda u'K$. Multiplication by $\lambda$ preserves the two
pairs of summands. For the base and every lift, the full equation
operator is a direct sum of two identical restricted $(I,C)$
operators. If their old solution dimension is $h_H$, then
\begin{equation}\label{eq:bip-line-splitting}
 h=2h_H.
\end{equation}
The quotient is isomorphic to the complementary restriction. It has
$p^{h_H}$ old point maps, each with an exact old lift through this
fixed complement. Translating that lift removes all affine forcing
at every $N$; an arbitrary affine forcing is not being assumed soluble.

The restricted full-uniform coefficient is exactly $p^{h_H}$.
Indeed its interaction dependencies have slopes $M,-M^T$.
The left map $I$ forces $M^T=-M$; the right map $C$ then forces
$CM=MC^T$. Cyclicity of $C$ kills $M$. The remaining constant
dependencies have dimension $h_H$, by the adjoint of the square
old two-map operator. With $Q_0=|K|$, the legal lattice index
$p^{kQ_0s}$ and compatible index $p^{kQ_0s-h_H}$ give coefficient
$p^{h_H}$ by the same exact normalization and Gaussian fraction.
All affine translates here come from literal exact old sections.

For each quotient point map, count only the restricted full-uniform
edge neighborhood inside its cosets. Its contribution is
$p^{h_H}+o(1)$, so one line pair contributes $p^{2h_H}+o(1)=p^h+o(1)$.
Using the entire restricted moment here would also count its point
subphases. Leakage outside these cosets contributes $o(1)$: the
weighted sparse-boundary bound is $O(N^{-1/2})$ times the bounded
restricted two-space moment. Its matching-forcing rank argument
controls the correlation with a large bulk kernel.
Summing all line pairs, point phases and original full uniformity
gives
\[
 (p^s+1)p^h+p^h+p^{h+s}=2(p^s+1)p^h,
\]
as required.

\subsection{Zero target and singleton-side coefficients}
Put $s_0=\dim\ker T=\dim\ker T^T$, so $h=2s_0$.
For scalar zero there are four subspace pairs:
$(0,0)$, $(\Fp,\Fp)$, $(\Fp,0)$ and $(0,\Fp)$.
The point family contributes $p^h+o(1)$ and full uniformity
contributes $p^{h+1}+o(1)$ by \eqref{eq:bip-uniform}.
For $(\Fp,0)$ the quotient point maps are
$a_R\in\ker T$, numbering $p^{s_0}$. Each has an exact section
with left labels zero and right labels fiber-constant. Its homogeneous
restriction is $X^Tz_L=0$.

We retain the full singleton-side prefactor. Left roots range over
$\Fp$ and right roots are zero. Before compatibility only right
old types have a nontrivial neighbor-sum congruence. These $m$
scalar congruences are independent and are all maximal periods,
by the same $d\ge3$ additive-character argument. Their full index
is $p^m$. Compatible integer directions are arbitrary left root
margins $n_i(a)$ of total zero; all edge tables are those margins
or their transposes. The moments $w_i=\sum_a a n_i(a)$ range
independently over $\Fp$ using label transfers between zero and
one. The restricted congruence map is $T^Tw$, of rank $m-s_0$.
Thus the compatible index is $p^{m-s_0}$.

There are $p^{dmN}$ legal star arrays in the iid normalization,
and the independent-star weighted normalizer after multiplication
by this factor is one. The full and compatible dimensions are
$m(d+1)(p-1)$ and $m(p-1)$; their difference cancels twice the
matching Stirling power. The remaining Gaussian exponent is
$-(p/2)\sum_{i\in\mathcal L}\|u_i\|^2$, with positive precision in every
nonzero direction. The relative density ratio $p^m$, compatible
fraction $p^{-(m-s_0)}$, and exact normalizer therefore give the
restricted uniform coefficient $p^{s_0}$.
The all-zero stars supply compatible affine origins at every $N$,
and the exact old section transports this calculation to each
quotient point map. This mixed family contributes
$p^{s_0}p^{s_0}+o(1)=p^h+o(1)$, as does the opposite family.

In each case choose narrow edge-table neighborhoods of the distinct
uniform-coset phase centers, together with small outside-coset radii.
There are finitely many centers. They can be chosen disjoint from
one another and from point centers. The sparse estimate makes
positive-defect leakage $o(1)$ using the bounded smaller-dimensional
two-space moment. The remaining compact feasible type region has
a strict entropy gap, so its exact factorial sum is exponentially
small. The same argument closes the nonzero cases above.
Adding the four zero families gives $(p+3)p^h+o(1)$, completing
the proof of Theorem~\ref{bip:thm:moments}.

\subsection{Vanishing excess probabilities and real transfer}
Fiber-constant pullback injects every old companion space into the
lift space over every field, even when $p\mid N$. For nonzero
irreducible $g$, both dimensions are multiples of $r=\deg g$.
Subtracting the deterministic baseline $p^h$ in Markov and using
\eqref{eq:bip-moments} gives
\begin{equation}\label{eq:bip-excess}
 \limsup_N\Prob\{\dim\ker g(A_N)>h\}\le
 \begin{cases}
 (p^r-1)^{-1},&g\ne t\text{ not even},\\
 (2p^s+1)/(p^{2s}-1),&g\text{ even of degree }2s,\\
 (p+2)/(p^2-1),&g=t.
 \end{cases}
\end{equation}
For the last denominator, the square blocks in
\eqref{eq:bip-square-blocks} give
$\dim\ker A_N=2\dim\ker X$ over every field. The old nullity
$h=2s_0$ is also even, so every excess is at least two. A
one-dimensional threshold would not give a vanishing zero-target
bound. For $p\ge3$ every bound in \eqref{eq:bip-excess} is at
most $2/(p-1)$.

To complete Theorem~\ref{lift:thm:all-base}, let $F$ be the squarefree
minimal polynomial of the real symmetric integer base $B$, and
fix an odd prime not dividing its nonzero discriminant. Then
$B\bmod p$ is semisimple. For an old irreducible integer factor
$f$ of degree $\ell$ and old exponent $t$, factor
$\overline f=\prod_jg_j$ into distinct irreducibles of degrees $r_j$.
Its old geometric dimensions are $\dim\ker g_j(B)=r_jt$.

A new real root of $f$ forces
$\dim_\Q\ker f(A_N)\ge\ell(t+1)$: the rational new characteristic
polynomial is divisible by $f$, and $A_N$ is real symmetric.
Reduction of the integer matrix $f(A_N)$ cannot increase rank.
Polynomial CRT gives
\[
 \ker\overline f(A_N\bmod p)
       =\bigoplus_j\ker g_j(A_N\bmod p)
\]
for any modular matrix, without semisimplicity of the lift.
Some factor must exceed its injected old baseline. Apply
\eqref{eq:bip-excess}, including its even increment at zero, and
sum over at most $\ell$ factors. This gives limiting collision
probability at most $2\ell/(p-1)$ for $f$, and at most
$2\deg F/(p-1)$ for all old roots. Arbitrarily large fixed good
primes tend to infinity only after the lift-size limit superior.
This proves the bipartite case and hence the all-base theorem.

This extension concerns simple regular bases with equal-size lift
fibers. It does not extend spectral disjointness to arbitrary
reversible degree frames: unequal numbers of old types on the two
bipartite sides can force new zero modes, as in Section~\ref{sec:scope}.
No simplicity of the new spectrum, numerical eigengap, or
growing-parameter uniformity is asserted.
